\section{Meaning}\label{sec:meaning}

What would a successful single proposal establish, and where would that
claim stop? The answer depends on the independent results in
\cref{sec:results}. The recorded training trajectory alone cannot establish
transfer, useful new composition or autonomous scientific discovery.
\TBD{Reconcile every conclusion below with the completed evaluation.}

\subsection{Learning to select and learning to invent}

A better task-dependent choice from a fixed program archive demonstrates
selection. A useful composition absent from the training archive is a stronger
result, and requires an output space that permits it. Token-by-token decoding does not distinguish the two:
report the admissible output space, the training archive and the relation
of each proposed program to that archive. \TBD{State which claim the evaluated
checkpoint and its outputs actually support.}

\subsection{Moving judgment across problems}

If learning improves first proposals on new families, it would support the
hypothesis that training evaluations can be retained as a transferable
proposal capability. It would not show that one proposal wins on every
problem. Test-time search may still improve a learned writer, and its value
should be compared with the training cost it takes to improve that writer.
\TBD{Use the measured search curves and full cost account to bound this claim.}

\subsection{The referee defines the boundary}

The framework optimizes what its referee measures. Repeated development
comparisons can overfit that referee, while a restricted language can hide
useful algorithms from the learner. Task representations, training families
and runtime limits further bound the result. A referee for a new scientific
domain would need its own validity and cost evidence; analogy to a game
supplies neither. \TBD{Report observed failures, distribution shifts and
remaining tests of the referee.}

\subsection{A bounded claim about self-improvement}

Improving an external algorithm, improving the proposer and improving the
rate at which the proposer learns are distinct claims. Evidence for the
first does not establish the others. People still choose the task, metric,
representation and resource limits. \TBD{State the strongest supported level
of improvement and leave untested extensions as future work.}

The contribution to assess is concrete: whether retained evaluation
experience improves the first executable proposal on a new problem, under
a declared budget and an independently scored test.
